\section[H. 朝复夕]{\myboxviii{      \BF{\TT{H}}} 朝复夕}

\textit{运行时间限制}   2 seconds

\textit{内存使用限制} 1 Gigabyte

\textit{读入文件} \qquad 标准输入

\textit{输出文件} \qquad 标准输出


\begin{quote}
   \textit{朝复夕，朝夕复。}
\end{quote}

\BF{本题和试题 \myboxvii{\BF{\TT{G}}} 朝复习没有关系。}

\subsection{题目描述}

小 T 准备出一些经典的问题 $\ldots$

对于给定一个长度为 $n$ 的排列 $p = \{p_1, p_2, \ldots, p_n\}$, 小 T 将进行一次操作 $(i, j)$：

\begin{itemize}
   \item 取一对下标 $1 \le i < j \le n$, 交换 $p$ 中的 $p_i$ 和 $p_j$.
\end{itemize}

对于所有不同的合法操作 $(i, j)$, 小 T 需要计算交换之后的排列所表示的置换的阶数之和。

由于答案可能很大，小 T 需要计算答案对 $998\,244\,353$ 取模的值。

请帮助小 T 计算这个问题。

试题的提示中附带有关于置换和排列的相关信息。

\subsection{输入格式}

第一行输入一个整数 $T$ ($1 \le T \le 10^5$), 表示测试数据的数量。

每组测试数据第一行输入一个整数 $n$ ($2 \le n \le 10^5$), 表示排列的长度。

随后一行输入 $n$ 个整数 $p_1$, $p_2$, $\ldots$, $p_n$ ($1 \le p_i \le n$), 表示所给定的排列。

保证所有测试数据输入中的 $\sum n \le 3 \times 10^5$, 且每组测试数据给定序列一定是一个合法的排列。

\subsection{输出格式}

对于每组测试数据输出一行一个整数，表示对于所有不同的合法操作 $(i, j)$, 交换之后的排列所表示的置换的阶数之和对 $998\,244\,353$ 取模的值。

\subsection{样例 1}

\begin{minipage}{\linewidth}
	\begin{minipage}[t]{0.49\linewidth}
		\begin{lstlisting}[title=\texttt{standard input}]
2
3
2 3 1
4
2 1 4 3
\end{lstlisting}
	\end{minipage}
	\begin{minipage}[t]{0.49\linewidth}
		\begin{lstlisting}[title=\texttt{standard output}]
6
20
\end{lstlisting}
	\end{minipage}
\end{minipage}


\subsection{样例 1 解释}

对于第 1 组测试数据，共有 $3$ 种不同的交换方式：
\begin{itemize}
   \item 交换 $p_1$ 和 $p_2$, 排列变为 $p = \{3, 2, 1\}$, 可以注意到 $p^2 = \text{id}_3$, 因为 $p(p(1)) = p(3) = 1$, $p(p(2)) = p(2) = 2$, $p(p(3)) = p(1) = 3$, 因此阶数为 $2$.
   \item 交换 $p_1$ 和 $p_3$, 排列变为 $p = \{1, 3, 2\}$, 可以注意到 $p^2 = \text{id}_3$, 阶数为 $2$.
   \item 交换 $p_2$ 和 $p_3$, 排列变为 $p = \{2, 1, 3\}$, 可以注意到 $p^2 = \text{id}_3$, 阶数为 $2$.
\end{itemize}

因此答案为 $6$.

\subsection{提示}

一个 $1$ 到 $n$ 的排列是一个长度为 $n$ 且其中每个 $1$ 到 $n$ 的整数都出现且仅出现一次的序列，这个排列的长度也为 $n$.

令 $S = \{1, 2, \ldots, n\}$, 一个 $1$ 到 $n$ 的置换是指从 $S$ 到 $S$ 的一个双射。


一个长度为 $n$ 的排列 $p = \{p_1, p_2, \ldots, p_n \}$可以表示一个置换 $p$, 也即 $p(i) = p_i$.

例如第 1 组样例的 $p = \{{\color[rgb]{0.36,0.65,0.59}2}, {\color[rgb]{0.75,0.41,0.46}3}, 1\}$ 表示的置换是 $p(n): \{1, 2, 3\} \rightarrow \{1, 2, 3\}$, 其中
\[
   p(1) = {\color[rgb]{0.36,0.65,0.59}2}, p(2) = {\color[rgb]{0.75,0.41,0.46}3}, p(3) = 1.
\]

长度为 $n$ 恒等置换 $\text{id}_n: \{1, 2, \ldots, n\} \rightarrow \{1, 2, \ldots, n\}$ 定义为 $\text{id}_n(i) = i$, 例如一个长度为 4 的恒等置换 $\text{id}_4$ 有：
\[
   \text{id}_4(1) = 1, \text{id}_4(2) = 2, \text{id}_4(3) = 3, \text{id}_4(4) = 4.
\]

两个置换 $p$ 和 $q$ 的复合 $p \circ q$ 表示的是 $p \circ q(i) = p(q(i))$ 这个置换。

长度为 $n$ 的置换的幂次 $p^m$ (其中 $m$ 是非负整数) 定义为
\[
   p^m = \begin{cases}
      p \circ p^{m - 1}, &m \ge 1;\\
      \text{id}_n, & m = 0.
   \end{cases}
\]

长度为 $n$ 的置换的\BF{阶数} $\text{ord}(p)$ 定义为最小的正整数 $k$, 满足 $p^{k} = \text{id}_n$.




